\section{Phương pháp}\label{sec:method}

Đề tài gồm một phân tích hồi cứu (RQ1, RQ2) và một thực nghiệm đối chứng (RQ3 đến RQ6). Phần hồi cứu
chỉ cho thống kê mô tả và là bối cảnh. Phần thực nghiệm là phần chính: mỗi nhiệm vụ được giao cho tác
tử theo hai cách làm, mỗi cách lặp ba lần, và chấm bằng test mà tác tử không thấy.

\subsection{Dữ liệu hồi cứu}

Kho mã ViPharmind được đọc ở hai mốc cố định: nhánh \texttt{main} tại commit
\texttt{3baae61d764f} và nhánh \texttt{camera-ready} tại commit \texttt{61d10db85f0f}. Tập commit phân
tích là hợp của hai mốc, mỗi commit tính một lần. Lấy hợp vì khâu nghiên cứu và viết bài nằm trên
\texttt{camera-ready}, còn sản phẩm nằm trên \texttt{main}. Tập này có \NumCommits{} commit, trong đó
\NumCommitsFeat{} commit \texttt{feat}, \NumCommitsFix{} \texttt{fix} và \NumCommitsRefactor{}
\texttt{refactor}. Không commit nào mang dòng ghi công AI (\texttt{Co-Authored-By}).

Mỗi commit được gán vào một khâu của vòng đời (nghiên cứu và đánh giá, dữ liệu, mã backend, mã
frontend, kiểm thử, triển khai, tài liệu, khác). Khâu của commit là khâu có nhiều dòng thay đổi nhất
trong commit, xác định theo luật đường dẫn file. Các file spec và plan được lấy lại từ lịch sử git vì
phần lớn đã bị xóa khỏi cây hiện tại; có \NumSpecPlanFiles{} file như vậy (\NumSpecFiles{} spec,
\NumPlanFiles{} plan).

\textit{Quy tắc nối commit với plan.} Một commit \texttt{feat} được coi là ``có plan đi trước'' khi nó
chạm ít nhất một file mã có đường dẫn được nêu trong một file plan đã được thêm vào trước commit đó
không quá 14 ngày. Để kiểm tra quy tắc, An lấy ngẫu nhiên \NumLinkChecked{} commit và đối chiếu tay
với plan. Trợ lý AI đọc trước từng cặp commit và plan rồi gợi ý một kết luận kèm lý do; An đọc và ghi
kết luận của mình. An đồng ý với gợi ý của AI ở cả \NumLinkAgreeAI{} commit. Độ chính xác của quy tắc
vì vậy là kết luận của An, nhưng không độc lập với AI.

\textit{Mức phải sửa lại (RQ2).} Với mỗi commit \texttt{feat}, đếm số commit \texttt{fix} hoặc
\texttt{refactor} chạm cùng file mã trong 14 ngày sau đó, rồi so sánh nhóm có plan đi trước với nhóm
không có.

\textit{Nhật ký phiên.} Nhật ký Claude Code còn giữ được gồm \NumSessionsMain{} phiên chính và
\NumSessionsSub{} nhật ký tác tử con. An đã xóa lịch sử chat giữa các phiên, nên phần này không đầy đủ
và chỉ dùng để mô tả.

\subsection{Bộ nhiệm vụ}

\textit{Nguồn và sàng lọc.} Ứng viên là mọi commit \texttt{feat} trên \texttt{main} chạm cả
\texttt{backend/src} và \texttt{backend/tests}. Mỗi nhiệm vụ gồm commit gốc, commit cha (mã nền), đề
bài và danh sách test ẩn. Test ẩn là các file test mà commit gốc thêm mới hoặc sửa. File thêm mới không
có trong thư mục làm việc; file bị sửa thì tác tử thấy bản cũ, và khi chấm bản của commit gốc ghi đè
lên. Bộ chạy áp tự động các tiêu chí 1, 4, 5 và 2 theo thứ tự đó (Bảng~\ref{tab:funnel}). Tiêu chí 2
được kiểm bằng cách chạy test thật ở cả commit cha và commit gốc.

\begin{table}[H]
\centering\small
\caption{Phễu sàng lọc nhiệm vụ}\label{tab:funnel}
\begin{tabular}{p{0.62\linewidth}rr}
\toprule
Bước & Còn lại & Bị loại \\
\midrule
Ứng viên: commit \texttt{feat} chạm cả \texttt{backend/src} và \texttt{backend/tests} & 110 &  \\
1. Không sửa file dùng chung của bộ test & 64 & 46 \\
4. Phần thay đổi ở \texttt{backend/src} từ 30 đến 400 dòng & 47 & 17 \\
5. Không chỉ đổi tên hay di chuyển file & 47 & 0 \\
2. Test ẩn hỏng ở commit cha, đạt ở commit gốc & 43 & 4 \\
\bottomrule
\end{tabular}

\end{table}

Tiêu chí 1 loại \NumRejectedCritOne{} commit, nhiều hơn con số khảo sát lúc thiết kế, vì bộ chạy dùng
đúng định nghĩa trong bản thiết kế: mọi file dưới \texttt{backend/tests/} không có tên dạng
\texttt{test\_*.py} đều là file dùng chung. Lý do của tiêu chí này: khi chấm, các file dùng chung được
khôi phục về bản gốc, việc đó sẽ ghi đè phần tác tử tự sửa. Tiêu chí 2 loại \NumRejectedCritTwo{}
commit; ba trong số đó có test mới đều thuộc loại \texttt{integration} (cần Docker), một commit làm
hỏng test cũ khi thêm test ẩn vào. Còn lại \NumEligible{} commit.

\textit{Tiêu chí 3 và duyệt đề.} Tiêu chí 3 đòi test ẩn kiểm tra hành vi qua giao diện mà đề bài nêu
được (đường dẫn API, chữ ký hàm, lớp công khai). Tiêu chí này cần đọc mã nên do người đánh giá. Trợ lý
AI đọc từng commit, đề xuất đạt hay không kèm lý do, và soạn nháp đề bài. Đề bài viết bằng tiếng Anh,
từ năm đến mười câu, nêu việc cần làm và giao diện phải giữ, không nêu cách cài đặt và không nhắc tới
test. An duyệt gộp toàn bộ file duyệt: chấp nhận nguyên văn mọi đề xuất tiêu chí 3 và mọi nháp đề bài,
không sửa từng mục. Ba mục được AI đề xuất đạt nhưng test hành vi chỉ thuộc loại \texttt{integration}
bị loại theo khuyến nghị của trợ lý. Một tác tử duyệt độc lập đã đối chiếu một mẫu đề bài với commit
gốc. Như vậy, đề bài là sản phẩm của AI, được người duyệt ở mức tổng thể. Mục~\ref{sec:discussion} bàn
hệ quả của việc này.

\textit{Lấy mẫu phân tầng.} \NumEligibleNoLayer{} commit không chạm lớp mã nào nên không phân tầng
được. \NumEligibleStratifiable{} commit còn lại được xếp thứ tự duyệt ngẫu nhiên với hạt giống cố định
\mbox{20261115}. Hạn mức theo lớp mã tỉ lệ với số commit đủ điều kiện của mỗi lớp (chia theo phần dư lớn nhất, mỗi lớp ít nhất một): \texttt{api}
\NumTasksApi{}, \texttt{application} \NumTasksApplication{}, \texttt{domain} \NumTasksDomain{},
\texttt{infrastructure} \NumTasksInfrastructure{}, tổng \NumTasks{} nhiệm vụ. Đi theo thứ tự duyệt,
mỗi commit được duyệt đạt sẽ vào lớp của nó nếu lớp còn hạn mức; hai commit đạt đầu tiên không vào được
hạn mức nào trở thành \NumPilotTasks{} nhiệm vụ thử, dùng để chạy thử bộ chạy. Phần thay đổi ở
\texttt{backend/src} của \NumTasks{} nhiệm vụ từ \NumTaskSrcLinesMin{} đến \NumTaskSrcLinesMax{} dòng,
trung vị \NumTaskSrcLinesMedian{} dòng. Danh sách đầy đủ ở Phụ lục~\ref{app:tasks}.

\subsection{Tầng tính năng (RQ6)}

Đề một commit nhỏ hơn nhiều so với một plan của ViPharmind, thường gồm khoảng mười phần việc. Đề nhỏ có
thể làm lợi ích của việc viết thiết kế trước không lộ ra, nên đề tài thêm một tầng nhiệm vụ cỡ một tính
năng. Gom nguyên một plan thành nhiệm vụ không làm được vì commit của các plan đan xen nhau. Thay vào
đó, phần lõi của từng plan backend được cắt thành các đoạn commit liền nhau, và bộ test thật được chạy
ở đầu và cuối mỗi đoạn. Ba đoạn được chọn là những đoạn mà test ẩn chỉ cần đề bài nêu ít tên và chạy
được độc lập. Mỗi đoạn có từ \NumFeatureCommitsMin{} đến \NumFeatureCommitsMax{} commit và từ
\NumFeatureSrcLinesMin{} đến \NumFeatureSrcLinesMax{} dòng thay đổi ở \texttt{backend/src}. Hai đoạn F09 và F04
bị loại vì đề bài sẽ phải nêu quá nhiều tên, gần như một bản thiết kế viết sẵn, làm nhiễu phép so sánh.
Đoạn F08 chứa commit của nhiệm vụ t03 và nhiệm vụ thử p02, nên hai tầng được phân tích riêng.

Mã nền của nhiệm vụ tính năng là commit \texttt{base} ngay trước đoạn. Test ẩn là mọi file \texttt{test\_*.py} thêm hoặc
sửa trong đoạn. Đề bài có thêm mục ``giao diện bắt buộc'', sinh tự động từ các tên mà test ẩn cần. Tầng này chỉ
chạy với hai mô hình Claude: \NumFeatureTasks{} nhiệm vụ, hai cách làm, ba lần lặp, hai mô hình, tổng
\NumFeatureRuns{} lượt. Kết quả chỉ được mô tả kèm khoảng tin cậy.

\subsection{Hai cách làm}

Hai cách làm nhận cùng đề bài, cùng mô hình, cùng quyền dùng công cụ, cùng giới hạn thời gian. Khác
biệt duy nhất là lời dặn về quy trình đặt trước đề bài, lưu thành hai file trong kho mã (nguyên văn ở
Phụ lục~\ref{app:prompts}). Lời dặn vibe coding chỉ có một câu: cài đặt thay đổi sau đây. Lời dặn theo
đặc tả yêu cầu ba giai đoạn trong cùng một phiên: viết \texttt{docs/design.md} (mục tiêu, phạm vi, giao
diện, cách làm, cách kiểm thử), viết \texttt{docs/plan.md} chia việc thành các phần trọn vẹn, rồi cài đặt
từng phần theo kiểu viết test trước. Lời dặn không quy định số phần hay độ nhỏ của phần, giống các plan
cũ của ViPharmind. Nó được viết dựa trên các spec và plan trong lịch sử dự án và không chép nguyên văn
skill nào.

Báo cáo gọi nhánh này là ``quy trình đặc tả tự động'': tác tử tự viết thiết kế và kế hoạch rồi cài đặt,
không có người duyệt. Quy trình An dùng khi làm ViPharmind có bước người duyệt spec và plan; thí nghiệm
không có bước đó. Nhánh đặc tả có người duyệt không được đưa vào vì spec gốc của ViPharmind được viết
cho những phần việc lớn hơn một commit và chỉ có ở một số nhiệm vụ. Lượt theo đặc tả mà không để lại cả
hai file thiết kế và kế hoạch được đánh dấu ``không tuân thủ quy trình'' và báo cáo riêng.

\subsection{Bộ chạy và bộ chấm}

Một đơn vị chạy là một tổ hợp (nhóm chạy, nhiệm vụ, cách làm, lần lặp). Bảng~\ref{tab:cohorts} liệt kê
các nhóm chạy. Mọi mô hình chạy ở mức effort hoặc mức suy luận medium. Phiên bản công cụ và mã mô hình
được ghi cho từng lượt.

\begin{table}[H]
\centering\small
\caption{Các nhóm chạy}\label{tab:cohorts}
\begin{tabularx}{\linewidth}{@{}llXl@{}}
\toprule
Đợt & Công cụ & Mô hình & Vai trò \\
\midrule
1 & Claude Code (\texttt{claude -p}) & Claude Sonnet 5.5 & RQ3, RQ5 \\
1 & Claude Code (\texttt{claude -p}) & Claude Opus 5.5 & RQ3, RQ5 \\
2 & Antigravity CLI (\texttt{agy -p}) & Gemini 3.8 Flash Medium & RQ4 \\
3 & Codex CLI (\texttt{codex exec}) & \texttt{gpt-6-luna} & RQ4 \\
\bottomrule
\end{tabularx}
\end{table}

Với tầng nhỏ, mỗi mô hình có \NumTasks{} nhiệm vụ, hai cách làm, ba lần lặp, tức \NumRunsPerModel{}
lượt; bốn mô hình qua ba đợt là \NumRunsSmallTotal{} lượt. Thứ tự các đơn vị được xáo trộn với hạt giống
cố định, để hai cách làm không rơi vào hai khung giờ khác nhau của dịch vụ. Lượt bị ngắt vì hết hạn mức
thuê bao không tính là thất bại và được chạy lại. Chi phí bằng tiền là số quy đổi từ lượng token thật
theo bảng giá API công khai, vì mọi lượt chạy bằng gói thuê bao.

Mỗi lượt chạy theo các bước sau. Bộ chạy xuất cây mã tại commit cha vào một thư mục mới, không kèm
lịch sử git, rồi tạo một commit nền duy nhất. Tác tử được gọi ở chế độ không tương tác với lời dặn và
đề bài, chạy tới khi tự dừng hoặc chạm giới hạn thời gian. Sau đó bộ chấm làm bốn việc. Thứ nhất, đo
chất lượng tĩnh (\texttt{ruff}, \texttt{pyrefly}, test kiến trúc) và chỉ số bảo trì (\texttt{radon}) trên
các file tác tử sửa. Thứ hai, đếm test tác tử tự viết và độ phủ của chúng trên phần mã tác tử thay đổi.
Thứ ba, dựng một bản sao không có file thiết kế và kế hoạch để dùng cho việc gán nhãn. Thứ tư, trên bản
sao đó, khôi phục mọi file test cũ về bản của commit cha, chép test ẩn của commit gốc vào, rồi chạy toàn
bộ test. Việc khôi phục test cũ ngăn tác tử ``đạt'' bằng cách sửa test. Bộ chấm cũng đánh dấu lượt nào
sửa \texttt{conftest.py} hoặc cấu hình \texttt{pytest}.

\subsection{Chống rò rỉ và điều kiện ngang nhau}

ViPharmind là repo công khai, và lịch sử của nó chứa lời giải. Bộ chạy chặn các đường rò rỉ sau. Cây mã
giao cho tác tử không có lịch sử git, nên tác tử không đọc được commit tương lai. Cây mã cũng bị xóa
mọi thư mục spec, plan cũ (\texttt{superpowers/}) và mọi file hướng dẫn hay cấu hình tác tử
(\texttt{CLAUDE.md}, \texttt{AGENTS.md}, \texttt{GEMINI.md}, \texttt{.claude/} và tương tự), vì plan cũ
có thể mô tả đúng tính năng cần làm; bộ chạy từ chối chạy nếu cây còn sót file nào như vậy. Công cụ truy
cập web của tác tử bị tắt. Mạng trong shell vẫn bật vì một trong ba công cụ không cô lập được nó; đây là
một hạn chế. Mỗi lượt dùng một thư mục cấu hình sạch (không plugin, không skill, không bộ nhớ, chỉ giữ
thông tin đăng nhập) và một thư mục làm việc tên ngẫu nhiên nằm ngoài kho mã nghiên cứu. Các lượt chạy
song song, nhưng hai lượt của cùng nhóm chạy và cùng nhiệm vụ không bao giờ chạy cùng lúc, để chúng
không thấy bài của nhau.

Ba công cụ có cơ chế khác nhau để tắt web, nạp cấu hình và đăng nhập. Bảng đối chiếu điều kiện chạy của
ba công cụ được lập sau lượt chạy thử: \pending{bảng đối chiếu điều kiện ba công cụ, từ báo cáo chạy
thử}.

\subsection{Chỉ số và phân tích}

Chỉ số chính là \emph{đạt nhiệm vụ}: mọi test ẩn đạt và không test cũ nào hỏng. Chỉ số phụ gồm tỉ lệ test
ẩn đạt, số test cũ bị hỏng (hồi quy), số lỗi \texttt{ruff} và \texttt{pyrefly}, kết quả test kiến trúc,
độ phức tạp chu trình và chỉ số bảo trì so với bản gốc, số test tác tử viết và độ phủ của chúng, số dòng
thêm và bớt, và chi phí (token vào, ra, chi phí quy đổi, thời gian chạy, số lượt trao đổi).

Phép so sánh luôn là hai cách làm trong cùng một nhóm chạy, theo cặp trên từng nhiệm vụ. Đề tài báo cáo
hiệu số tỉ lệ đạt kèm khoảng tin cậy bootstrap lấy mẫu lại theo nhiệm vụ, và độ dao động giữa ba lần
lặp. Với số nhiệm vụ nhỏ, kết luận dựa vào cỡ hiệu ứng và khoảng tin cậy, không dựa vào ngưỡng p. Ba công
cụ dùng mô hình khác hãng và khác hạng, nên đề tài không so sánh công cụ với công cụ. Ngoại lệ là Sonnet
với Opus trong Claude Code (RQ5), cùng công cụ và cùng hãng. Đợt 1 được phân tích như thiết kế
$2\times 2$ (cách làm $\times$ mô hình): bốn tỉ lệ đạt, hiệu số giữa hai cách làm ở từng mô hình, và hiệu
của hai hiệu số đó kèm khoảng tin cậy. Theo kích thước nhiệm vụ, \NumTasks{} nhiệm vụ được chia ba nhóm
theo số dòng của bản vá gốc; phân tích này là khám phá.

Hai thành viên Lê Xuân Lợi và Trần Ngọc Bảo gán nhãn độc lập nguyên nhân thất bại cho mọi lượt không đạt
của đợt 1, theo bảng nhãn có định nghĩa và ví dụ. Họ nhận log test và bản vá đã bỏ file thiết kế, kế
hoạch, nên không biết lượt đó thuộc cách làm nào. Độ đồng thuận được báo cáo bằng hệ số kappa của Cohen;
chỗ lệch nhau do An phân xử. An cũng đọc một mẫu ngẫu nhiên lượt đạt, đối chiếu với bản vá gốc, để tìm
trường hợp đạt test nhưng sai.

\subsection{Chốt trước khi chạy}

Danh sách nhiệm vụ, đề bài, lời dặn, giả thuyết, chỉ số và kế hoạch phân tích được commit và gắn thẻ git
\texttt{prereg-v1} trước lượt chạy chính đầu tiên. Lượt chạy trên các nhiệm vụ thử được phép trước thẻ
này, và giới hạn thời gian của mỗi lượt được chốt từ đó. Mọi thay đổi sau thẻ được ghi lại kèm lý do
trong mục~\ref{sec:results}.

\subsection{Vai trò của AI trong chính nghiên cứu này}

Theo hướng dẫn của \cite{baltes_2025}, đề tài khai báo mọi việc có AI tham gia. AI là đối tượng được
đo trong thực nghiệm. AI đã được dùng để làm ViPharmind. Trợ lý AI soạn nháp đề bài và đề xuất đánh giá
tiêu chí 3; An duyệt gộp như mô tả ở trên. Trợ lý AI gợi ý kết luận khi An kiểm tra quy tắc nối commit
với plan. Trợ lý AI viết phần lớn mã của bộ chạy, bộ chấm, phân tích hồi cứu và script sinh số liệu cho
báo cáo; An duyệt thiết kế của từng phần. Trợ lý AI cũng viết nháp báo cáo này; mọi con số trong văn bản
được sinh bằng máy từ file kết quả, và mọi trích dẫn đã được mở và kiểm tra ở trang gốc.
